% GENERATED FILE. Edit canonical lessons, then run npm run generate:books.
% canonical-source-hash: fnv1a64:7c1fbfe75a77a54f
% canonical-lessons: JA-C81-atari, JA-C81-fumoto, JA-C81-tokai, JA-C81-aida, JA-C81-chihou, JA-R81-first-pass-times-of-day-and-dates, JA-R81-second-pass-places-and-directions

\chapter{Area around, Foot of a mountain, City, Between, Region}
\label{ch:ja-area-around-foot-of-a-mountain-city-between-region}
\hlchaptermodality{81}

\begin{chapteropening}
\textbf{By the end of this chapter:} \emph{I can say the area around, the foot of a mountain, the city, between and a region.}

Complete the last lesson of chapter 81: i can say the area around, the foot of a mountain, the city, between and a region.
\end{chapteropening}

\section[\texorpdfstring{\ja{あたり} (atari)}{あたり (atari)}]{\ja{あたり} (atari) — the area around}

\label{lesson:JA-C81-atari}



\begin{quote}
Before the new one: say the Japanese for a shortcut, then the Japanese for the surroundings.
\end{quote}

\subsection*{You'll want to know: \ja{あたり}}
\textbf{\ja{あたり}} — \emph{atari} — \textquotedblleft{}the area around\textquotedblright{}.

\begin{grammarlens}[title={What you've built}]
One more word for this chapter.
\end{grammarlens}

\subsection*{Your turn}
\begin{itemize}
\raggedright
  \item \emph{Say it:} \emph{atari}
  \item \emph{Say it:} \emph{atari}, once more
  \item \emph{Say it:} say \emph{atari}
  \item \emph{Recall:} say the Japanese for a shortcut, then the Japanese for the surroundings, then say \emph{atari} again
\end{itemize}

\begin{tcolorbox}[breakable,colback=teal!4,colframe=teal!35!black,title={Before you move on}]
What does \ja{あたり} mean? (\textquotedblleft{}The area around\textquotedblright{}.) Say it once more.
\end{tcolorbox}
\section[\texorpdfstring{\ja{ふもと} (fumoto)}{ふもと (fumoto)}]{\ja{ふもと} (fumoto) — the foot of a mountain}

\label{lesson:JA-C81-fumoto}



\begin{quote}
Before the new one: say the Japanese for the surroundings, then the Japanese for the area around.
\end{quote}

\subsection*{You'll want to know: \ja{ふもと}}
\textbf{\ja{ふもと}} — \emph{fumoto} — \textquotedblleft{}the foot of a mountain\textquotedblright{}.

\begin{grammarlens}[title={What you've built}]
One more word for this chapter.
\end{grammarlens}

\subsection*{Your turn}
\begin{itemize}
\raggedright
  \item \emph{Say it:} \emph{fumoto}
  \item \emph{Say it:} \emph{fumoto}, once more
  \item \emph{Say it:} say \emph{fumoto}
  \item \emph{Recall:} say the Japanese for the surroundings, then the Japanese for the area around, then say \emph{fumoto} again
\end{itemize}

\begin{tcolorbox}[breakable,colback=teal!4,colframe=teal!35!black,title={Before you move on}]
What does \ja{ふもと} mean? (\textquotedblleft{}The foot of a mountain\textquotedblright{}.) Say it once more.
\end{tcolorbox}
\section[\texorpdfstring{\ja{とかい} (tokai)}{とかい (tokai)}]{\ja{とかい} (tokai) — the city}

\label{lesson:JA-C81-tokai}



\begin{quote}
Before the new one: say the Japanese for the area around, then the Japanese for the foot of a mountain.
\end{quote}

\subsection*{You'll want to know: \ja{とかい}}
\textbf{\ja{とかい}} — \emph{tokai} — \textquotedblleft{}the city\textquotedblright{}.

\begin{grammarlens}[title={What you've built}]
One more word for this chapter.
\end{grammarlens}

\subsection*{Your turn}
\begin{itemize}
\raggedright
  \item \emph{Say it:} \emph{tokai}
  \item \emph{Say it:} \emph{tokai}, once more
  \item \emph{Say it:} say \emph{tokai}
  \item \emph{Recall:} say the Japanese for the area around, then the Japanese for the foot of a mountain, then say \emph{tokai} again
\end{itemize}

\begin{tcolorbox}[breakable,colback=teal!4,colframe=teal!35!black,title={Before you move on}]
What does \ja{とかい} mean? (\textquotedblleft{}The city\textquotedblright{}.) Say it once more.
\end{tcolorbox}
\section[\texorpdfstring{\ja{あいだ} (aida)}{あいだ (aida)}]{\ja{あいだ} (aida) — between}

\label{lesson:JA-C81-aida}



\begin{quote}
Before the new one: say the Japanese for the foot of a mountain, then the Japanese for the city.
\end{quote}

\subsection*{You'll want to know: \ja{あいだ}}
\textbf{\ja{あいだ}} — \emph{aida} — \textquotedblleft{}between\textquotedblright{}.

\begin{grammarlens}[title={What you've built}]
One more word for this chapter.
\end{grammarlens}

\subsection*{Your turn}
\begin{itemize}
\raggedright
  \item \emph{Say it:} \emph{aida}
  \item \emph{Say it:} \emph{aida}, once more
  \item \emph{Say it:} say \emph{aida}
  \item \emph{Recall:} say the Japanese for the foot of a mountain, then the Japanese for the city, then say \emph{aida} again
\end{itemize}

\begin{tcolorbox}[breakable,colback=teal!4,colframe=teal!35!black,title={Before you move on}]
What does \ja{あいだ} mean? (\textquotedblleft{}Between\textquotedblright{}.) Say it once more.
\end{tcolorbox}
\section[\texorpdfstring{\ja{ちほう} (chihou)}{ちほう (chihou)}]{\ja{ちほう} (chihou) — a region}

\label{lesson:JA-C81-chihou}



\begin{quote}
Before the new one: say the Japanese for the city, then the Japanese for between.
\end{quote}

\subsection*{You'll want to know: \ja{ちほう}}
\textbf{\ja{ちほう}} — \emph{chihou} — \textquotedblleft{}a region\textquotedblright{}.

\begin{grammarlens}[title={What you've built}]
One more word for this chapter.
\end{grammarlens}

\subsection*{Your turn}
\begin{itemize}
\raggedright
  \item \emph{Say it:} \emph{chihou}
  \item \emph{Say it:} \emph{chihou}, once more
  \item \emph{Say it:} say \emph{chihou}
  \item \emph{Recall:} say the Japanese for the city, then the Japanese for between, then say \emph{chihou} again
\end{itemize}

\begin{tcolorbox}[breakable,colback=teal!4,colframe=teal!35!black,title={Before you move on}]
What does \ja{ちほう} mean? (\textquotedblleft{}A region\textquotedblright{}.) Say it once more.
\end{tcolorbox}
\section[(dialogue)]{Times of day and dates — a first pass}

\label{lesson:JA-R81-first-pass-times-of-day-and-dates}



\begin{quote}
Say the words for between and a region.
\end{quote}

\subsection*{Your turn}
Say these aloud:

\begin{itemize}
\raggedright
  \item the afternoon, the middle of the night, early morning, the evening meal, the morning sun
  \item sunset, tonight, the day after tomorrow, the day before yesterday, every morning
  \item every week, the weekend, the end, on the way, half past
  \item every year, this year, this week, the year before last, oversleeping
  \item midnight, a whole week, one whole day, half a day, the first of the month
\end{itemize}

\begin{tcolorbox}[breakable,colback=teal!4,colframe=teal!35!black,title={Before you move on}]
The afternoon? (\textbf{\ja{ごご}}, \emph{gogo}.) The middle of the night? (\textbf{\ja{よなか}}, \emph{yonaka}.) Early morning? (\textbf{\ja{あさがた}}, \emph{asagata}.) The evening meal? (\textbf{\ja{ゆうはん}}, \emph{yuuhan}.) The morning sun? (\textbf{\ja{あさひ}}, \emph{asahi}.) Sunset? (\textbf{\ja{ひのいり}}, \emph{hinoiri}.) Tonight? (\textbf{\ja{こんや}}, \emph{konya}.) The day after tomorrow? (\textbf{\ja{あさって}}, \emph{asatte}.) The day before yesterday? (\textbf{\ja{おととい}}, \emph{ototoi}.) Every morning? (\textbf{\ja{まいあさ}}, \emph{maiasa}.) Every week? (\textbf{\ja{まいしゅう}}, \emph{maishuu}.) The weekend? (\textbf{\ja{しゅうまつ}}, \emph{shuumatsu}.) The end? (\textbf{\ja{おわり}}, \emph{owari}.) On the way? (\textbf{\ja{とちゅう}}, \emph{tochuu}.) Half past? (\textbf{\ja{はん}}, \emph{han}.) Every year? (\textbf{\ja{まいとし}}, \emph{maitoshi}.) This year? (\textbf{\ja{ことし}}, \emph{kotoshi}.) This week? (\textbf{\ja{こんしゅう}}, \emph{konshuu}.) The year before last? (\textbf{\ja{おととし}}, \emph{ototoshi}.) Oversleeping? (\textbf{\ja{あさねぼう}}, \emph{asanebou}.) Midnight? (\textbf{\ja{まよなか}}, \emph{mayonaka}.) A whole week? (\textbf{\ja{いっしゅうかん}}, \emph{isshuukan}.) One whole day? (\textbf{\ja{いちにち}}, \emph{ichinichi}.) Half a day? (\textbf{\ja{はんにち}}, \emph{hannichi}.) The first of the month? (\textbf{\ja{ついたち}}, \emph{tsuitachi}.)
\end{tcolorbox}
\section[(dialogue)]{Places and directions — a second pass}

\label{lesson:JA-R81-second-pass-places-and-directions}



\begin{quote}
Say the words for the second of the month and the third of the month.
\end{quote}

\subsection*{Your turn}
Say these aloud:

\begin{itemize}
\raggedright
  \item the second of the month, the third of the month, the fourth of the month, the fifth of the month; some day, the sixth of the month
  \item where, nearby, behind, beside, the very middle
  \item the other side, this side of, a corner, stairs, a corridor
  \item a crossroads, a traffic light, a direction, a shortcut, the surroundings
  \item the area around, the foot of a mountain, the city, between, a region
\end{itemize}

\begin{tcolorbox}[breakable,colback=teal!4,colframe=teal!35!black,title={Before you move on}]
The second of the month? (\textbf{\ja{ふつか}}, \emph{futsuka}.) The third of the month? (\textbf{\ja{みっか}}, \emph{mikka}.) The fourth of the month? (\textbf{\ja{よっか}}, \emph{yokka}.) The fifth of the month; some day? (\textbf{\ja{いつか}}, \emph{itsuka}.) The sixth of the month? (\textbf{\ja{むいか}}, \emph{muika}.) Where? (\textbf{\ja{どこ}}, \emph{doko}.) Nearby? (\textbf{\ja{ちかく}}, \emph{chikaku}.) Behind? (\textbf{\ja{うしろ}}, \emph{ushiro}.) Beside? (\textbf{\ja{よこ}}, \emph{yoko}.) The very middle? (\textbf{\ja{まんなか}}, \emph{mannaka}.) The other side? (\textbf{\ja{むこう}}, \emph{mukou}.) This side of? (\textbf{\ja{てまえ}}, \emph{temae}.) A corner? (\textbf{\ja{かど}}, \emph{kado}.) Stairs? (\textbf{\ja{かいだん}}, \emph{kaidan}.) A corridor? (\textbf{\ja{ろうか}}, \emph{rouka}.) A crossroads? (\textbf{\ja{こうさてん}}, \emph{kousaten}.) A traffic light? (\textbf{\ja{しんごう}}, \emph{shingou}.) A direction? (\textbf{\ja{ほうがく}}, \emph{hougaku}.) A shortcut? (\textbf{\ja{ちかみち}}, \emph{chikamichi}.) The surroundings? (\textbf{\ja{まわり}}, \emph{mawari}.) The area around? (\textbf{\ja{あたり}}, \emph{atari}.) The foot of a mountain? (\textbf{\ja{ふもと}}, \emph{fumoto}.) The city? (\textbf{\ja{とかい}}, \emph{tokai}.) Between? (\textbf{\ja{あいだ}}, \emph{aida}.) A region? (\textbf{\ja{ちほう}}, \emph{chihou}.)
\end{tcolorbox}
